\section*{章末练习}
\begingroup\small
\subsection*{口述概念}
\begin{enumerate}[nosep,leftmargin=*]
  \item 区分序列收敛、Cauchy 性与空间完备性。为什么“迭代内部越来越接近”还不能保证存在允许状态空间内的极限？
  \item 区分函数极限、连续性与一致连续性，并解释函数极限为什么要求目标点是聚点。
  \item 区分函数序列的逐点收敛与一致收敛。为什么反例中的见证点可以随 $n$ 改变？
\end{enumerate}

\subsection*{笔试推导}
\begin{enumerate}[nosep,leftmargin=*]
  \item 使用三角不等式证明度量空间中的序列极限唯一，并标明分离性在哪一步使用。
  \item 从压缩条件推导相邻差、Cauchy 几何级数界、先验误差界和后验误差界；指出完备性与连续性分别在哪里进入。
  \item 证明连续函数的一致极限仍连续；再用该结论或移动见证证明 $f_n(x)=x^n$ 在 $[0,1]$ 上不一致收敛。
\end{enumerate}

\subsection*{数值编程}
\begin{enumerate}[leftmargin=*]
  \item 迭代 $x_{n+1}=0.02+0.8x_n,x_0=0$，显示数值、步长与误差。
  \item 保持固定点为 0.1，把压缩率改为 0.99，比较所需步数。
  \item 对 $x^n$ 比较固定点 $x=1/2$ 与移动点 $x_n=2^{-1/n}$。
\end{enumerate}

\subsection*{研究判断}
\begin{enumerate}[nosep,leftmargin=*]
  \item 优化器连续 20 步满足 $\lVert x_{n+1}-x_n\rVert_2<10^{-6}$。这能否证明收敛到最优解？列出仍缺失的数学与实现证据。
  \item 一个算法的状态被限制为有理数，却逼近 $\sqrt2$。继续增加迭代次数能否解决“不返回合法极限”的问题？给出模型层修复。
  \item 神经网络逼近在一万个验证网格点上误差都小于 $10^{-3}$。研究员据此声称连续输入域上一致误差小于 $10^{-3}$；评价该结论并设计更强证据。
\end{enumerate}
\endgroup

\subsection*{基础衔接：独立计算}
在 $[0,1/2]$ 上证明 $x^n$ 一致趋于零，并给出误差不超过 $10^{-3}$ 所需的最小整数 $n$。
